\centerline{\bf \Huge Математические пострелушки.}

\begin{flushright}
        \scriptsize{Неважно, возможно это или нет. Я это сделаю, потому что хочу сделать. \\
        Монки Д. Луффи}
\end{flushright}


\section{Оружейная 1}

\problem{1;1} На одну чашку весов положили круг сыра, а на другую — такого же круга и еще килограммовую гирю. Установилось равновесие. Сколько весит круг сыра?

\problem{2;2}Расставьте знаки арифметических действий и скобки, чтобы получились верные равенства: \\
$ \textbf{а}) 4 \quad 4 \quad 4 \quad 4 =5; \\  
            \textbf{б}) 4 \quad 4 \quad 4 \quad 4 =17;\\ 
            \textbf{в}) 5 \quad 5 \quad 5 \quad  5 = 30$ 

\problem{3;3} Для нумерации страниц в книге потребовалось 2 322 цифры. Сколько страниц в книге, если нумерация начинается с первой страницы.)

\problem {4;2} Разрежьте фигуру (см. рисунок) по линиям сетки на четыре равные фигуры.
\begin{center}
\includegraphics[width=0.2\textwidth]{Летние Сборы/Сборы Элиста 2024/ВВ/пострелушки_2.png}
\end{center}

\problem{Ыы} Найдите сумму чисел $1 + 2 + 3 + \dots + 2049 $

\problem{5;1}
Во время игры в шахматы у Эй Арнольда в какой-то момент оказалось на доске в два раза меньше фигур, чем у Медвежонка, при этом их было в пять раз меньше чем свободных клеток на доске. Сколько фигур Медвежонка было съедено к этому моменту?



\newpage

\section{Оружейная 2}

\problem{1;1}
Придумайте трехзначное число, у которого с любым из чисел 543, 142 и 562 совпадает один из разрядов, а два других не совпадают.

\problem{2;2}
Известно, что 9 стаканов отборного кваса стоят дешевле 10 рублей, а 10 стаканов отборного кваса - дороже 11 рублей. Сколько стоит стакан отборного кваса? 

\problem{3;1} Бобер пилит бревно. Он сделали 10 распилов. Сколько получилось кусков?

\problem{ыЫ}  
Найдите наименьшее натуральное число кратное 100, сумма цифр которого равна 100.

\problem{4;2} Запишите число, являющееся суммой 13 тысяч, 12 сотен и 11 единиц.


\problem{5;3}
Заполните свободные клетки "шестиугольника" (см. рисунок) целыми числами от 1 до 19 так, чтобы во всех вертикальных и диагональных рядах сумма чисел, стоящих в одном ряду, была бы одна и та же.
\begin{center}
\includegraphics[width=0.2\textwidth]{Летние Сборы/Сборы Элиста 2024/ВВ/Пострелушки_1.png}
\end{center}

\newpage


\section{Оружейная 3}

\problem{1;2}
Решите ребус:  $LEGO + LES = 2011.$

\problem{Ы?} Угол при вершине журавлиного клина равен 20°. 
Как изменится величина этого угла при рассматривании журавлей в бинокль с троекратным увеличением?

\problem{2;1} На какое максимальное число кусков можно разделить круглый блинчик при помощи трех прямолинейных разрезов?

\problem{3;2}
Составитель с Вячеславом ходил в тир. Уговор был такой: Составитель делает 5 выстрелов и за каждое попадание в цель получает право сделать ещё два выстрела. Всего Составитель сделал 17 выстрелов. Сколько раз Составитель попал в цель?

\problem{4;3}
В комнате стоят 20 стульев двух цветов: синего и красного. На каждый из стульев сел либо рыцарь, либо лжец. Рыцари всегда говорят правду, лжецы всегда лгут. Каждый из сидящих заявил, что он сидит на синем стуле. Затем они как-то пересели, после чего половина из сидящих сказали, что сидят на синих стульях, а остальные сказали, что сидят на красных. Сколько рыцарей теперь сидит на красных стульях? 

\problem{5;1} 
У 28 человек ЭЛМШ «Ы» класса на собрание пришли папы и мамы. Мам было - 24, пап - 18. У скольких учеников на собрание пришли одновременно и папа и мама?
